\begin{lemma}[Reset between consecutive colour rounds]
Let $B$ be a blow-up instance with finite source and host vertex sets,
candidacy system $A$, and reduced graph $R$ on $[r]$.  Fix a colouring
$c:[r]\to[q]$, sequences $\varepsilon,p:\mathbb N\to[0,\infty]$,
$D\in\mathbb N$, $C\in[0,\infty]$, and $n,t\in\mathbb N$.
Let $B_t=\{j:c(j)=t\}$, $X$ be a finite source-vertex set, and $\mu$ a
specified law on compatible partial embeddings with domain $X$.
Assume $t+1<q$, $\varepsilon_t\leq\varepsilon_{t+1}$, and the completed-class
invariant $\mathcal I(t,B_t;X,\mu)$ of \noderef{ColourRoundInvariant}, with
parameters $(\varepsilon,p,D,C,n)$.  Then the same domain and the same law
satisfy
\[
 \mathcal I(t+1,\varnothing;X,\mu).
\]
\end{lemma}
\begin{proof}
We verify every clause of \noderef{ColourRoundInvariant}.  This is the
round-boundary argument of \texttt{paper/r2-colour-round-addendum.tex},
lines 241--271 and 333--341, restricted to boundaries with $t+1<q$.
Put $P=\{k:c(k)<t\}\cup B_t$ and $P'=\{k:c(k)<t+1\}$.
For each cluster $k$, natural-number arithmetic gives
\[
 k\in P\ \Longleftrightarrow\ c(k)<t\text{ or }c(k)=t
       \ \Longleftrightarrow\ c(k)<t+1
       \ \Longleftrightarrow\ k\in P'.
\]
Thus $P=P'$, including when $B_t$ is empty.  Applying this equality to
the source-cluster index of every source vertex proves that the old domain
clause is exactly the new domain clause, with the same $X$.
The degree bound and square-colouring clauses are unchanged, the new touched
set is empty and hence contained in the next colour class, and the new
round is legal by $t+1<q$.  The law and its restriction maps are unchanged,
so its vertex-spread clause, interpreted by \noderef{VertexSpreadOn},
is the old clause without any change of probability law.

Fix $f\in\operatorname{supp}\mu$ and a cluster $j$ for which the new witness
clause is required, namely $c(j)\geq t+1$.  Then $j\notin B_t$ and
$c(j)\geq t$, so the old invariant supplies its witness.  Its old and new
exposure finsets are equal, not just equicardinal:
\[
 \{k\in P:R.Adj(j,k)\}=\{k\in P':R.Adj(j,k)\}.
\]
Let their common cardinality be $x_j$.  The old bound $x_j\leq D$ therefore
remains valid, and the density is the very same $p_{x_j}$.
The new touched-neighbour finset is empty, so its cardinality is zero and
in particular at most one.

The partial embedding $f$ itself is unchanged, in the same type with domain
$X$, as in \noderef{PartialCompatibleEmbedding}.  Consequently $f(X)$,
the unused source part of cluster $j$, and its unused host part are all
unchanged.  Unfolding \noderef{UpdatedCandidacy} shows that every allowed
pair is unchanged: it tests the same domain exclusion, image exclusion,
original allowed pair, and adjacency to every $f(x')$ with $x'\in X$.
Hence \noderef{UpdatedCandidacyCluster} gives the identical updated graph
of cluster $j$ before and after the reset.  The witness graph $W$ supplied
by the old invariant is therefore still contained in that same graph on
the same two parts $L,Q$.

If $N_R(j)\cap B_t$ is nonempty, this old witness already has error
$\varepsilon_{t+1}$, exactly the new error because the new touched set is
empty.  Keep its super-regularity assertion unchanged.  If the intersection
is empty, the old witness has error $\varepsilon_t$.  Write
$a=\varepsilon_t$, $b=\varepsilon_{t+1}$, and $d=p_{x_j}$, so $a\leq b$.
We check directly that its $(a,d)$-super-regularity implies
$(b,d)$-super-regularity.
By \noderef{RegularPair}, for subsets $U\subseteq L$ and $V\subseteq Q$
with $|U|>b|L|$ and $|V|>b|Q|$, monotonicity of multiplication gives
$|U|>a|L|$ and $|V|>a|Q|$.  The old regularity supplies
$d-a\leq d_W(U,V)\leq d+a$.
Truncated subtraction is antitone in its second argument and addition is
monotone, so
\[
 d-b\leq d-a\leq d_W(U,V)\leq d+a\leq d+b.
\]
This proves the new regularity, with the strict thresholds as required by
the definition.  For each $u\in L$, the old left degree bounds of
\noderef{SuperRegularPair} give
\[
 (d-b)|Q|\leq(d-a)|Q|\leq\deg_W(u,Q)
       \leq(d+a)|Q|\leq(d+b)|Q|.
\]
For each $v\in Q$, the identical monotonicities give the right degree bounds
\[
 (d-b)|L|\leq(d-a)|L|\leq\deg_W(v,L)
       \leq(d+a)|L|\leq(d+b)|L|.
\]
These inequalities use the extended nonnegative-real arithmetic of the
definitions and remain valid if a part is empty or a parameter is infinite.
Together they prove the new super-regularity of the unchanged $W$.
Thus only previously untouched witnesses have had their error weakened;
all witness, domain, colouring, degree, exposure, and spread clauses hold
with the same $X$ and $\mu$.  This proves the claimed next-round invariant.
There is no assertion at round $q$, since the hypothesis excludes it.
\end{proof}
